\documentclass[12pt,a4paper]{ctexbook}
\usepackage{geometry}
\geometry{margin=2.2cm}
\usepackage{setspace}
\setstretch{1.35}
\usepackage{hyperref}
\title{花间集·huajianji-2-juan}
\author{古典文献整理}
\date{}
\begin{document}
\maketitle
\tableofcontents
\newpage
\section{清平乐·上陽春晚}
\textbf{温庭筠}\par\vspace{0.5em}
上陽春晚，宫女愁蛾浅。\par
新岁清平思同辇，争那长安路远。\par
凤帐鸳被徒熏，寂寞花锁千门。\par
竞把黄金买赋，为妾将上明君。\par

\section{清平乐·洛陽愁绝}
\textbf{温庭筠}\par\vspace{0.5em}
洛陽愁绝，杨柳花飘雪。\par
终日行人恣攀折，桥下流水呜咽。\par
上马争劝离觞，南浦莺声断肠。\par
愁杀平原年少，回首挥泪千行。\par

\section{遐方怨·凭绣槛}
\textbf{温庭筠}\par\vspace{0.5em}
凭绣槛，解罗帷。\par
未得君书，肠断，潇湘春雁飞。\par
不知征马几时归，海棠花谢也，雨霏霏。\par

\section{遐方怨·花半坼}
\textbf{温庭筠}\par\vspace{0.5em}
花半坼，雨初晴。\par
未卷珠帘，梦残，惆怅闻晓莺。\par
宿妆眉浅粉山横，约鬟鸾镜里，绣罗轻。\par

\section{诉衷情·莺语花舞春昼午}
\textbf{温庭筠}\par\vspace{0.5em}
莺语花舞春昼午，雨霏微。\par
金带枕，宫锦，凤凰帷。\par
柳弱蝶交飞，依依。\par
辽陽音信稀，梦中归。\par

\section{思帝乡·花花}
\textbf{温庭筠}\par\vspace{0.5em}
花花，满枝红似霞。\par
罗袖画帘肠断，卓香车。\par
回面共人闲语，战篦金凤斜。\par
惟有阮郎春尽，不归家。\par

\section{梦江南·千万恨}
\textbf{温庭筠}\par\vspace{0.5em}
千万恨，恨极在天涯。\par
山月不知心里事，水风空落眼前花，摇曳碧云斜。\par

\section{梦江南·梳洗罢}
\textbf{温庭筠}\par\vspace{0.5em}
梳洗罢，独倚望江楼。\par
过尽千帆皆不是，斜晖脉脉水悠悠，肠断白苹洲。\par

\section{河传·江畔}
\textbf{温庭筠}\par\vspace{0.5em}
江畔，相唤。\par
晓妆鲜，仙景个女采莲。\par
请君莫向那岸边，少年，好花新满船。\par
红袖摇曳逐风暖，垂玉腕，肠向柳丝断。\par
浦南归，浦北归，莫知，晚来人已稀。\par

\section{河传·湖上}
\textbf{温庭筠}\par\vspace{0.5em}
湖上，闲望。\par
雨萧萧，烟浦花桥路遥。\par
谢娘翠蛾愁不消，终朝，梦魂迷晚潮。\par
荡子天涯归棹远，春已晚，莺语空肠断。\par
若耶溪，溪水西，柳堤，不闻郎马嘶。\par

\section{河传·同伴}
\textbf{温庭筠}\par\vspace{0.5em}
同伴，相唤。\par
杏花稀，梦里每愁依违。\par
仙客一去燕已飞，不归，泪痕空满衣。\par
天际云鸟引晴远，春已晚，烟霭渡南苑。\par
雪梅香，柳带长，小娘，转令人意伤。\par

\section{番女怨·万枝香雪开已遍}
\textbf{温庭筠}\par\vspace{0.5em}
万枝香雪开已遍，细雨双燕。\par
钿蝉筝，金雀扇，画梁相见。\par
雁门消息不归来，又飞回。\par

\section{番女怨·碛南沙上惊雁起}
\textbf{温庭筠}\par\vspace{0.5em}
碛南沙上惊雁起，飞雪千里。\par
玉连环，金镞箭，年年征战。\par
画楼离恨锦屏空，杏花红。\par

\section{荷叶杯·一点露珠凝冷}
\textbf{温庭筠}\par\vspace{0.5em}
一点露珠凝冷，波影。满池塘。\par
绿茎红艳两相乱，肠断。水风凉。\par

\section{荷叶杯·镜水夜来秋月}
\textbf{温庭筠}\par\vspace{0.5em}
镜水夜来秋月，如雪。采莲时。\par
小娘红粉对寒浪，惆怅。正思惟。\par

\section{荷叶杯·楚女欲归南浦}
\textbf{温庭筠}\par\vspace{0.5em}
楚女欲归南浦，朝雨。湿愁红。\par
小船摇漾入花里，波起。隔西风。\par

\section{天仙子·晴野鹭鸶飞一只}
\textbf{皇甫松}\par\vspace{0.5em}
晴野鹭鸶飞一只，水荭花发秋江碧。\par
刘郎此日别天仙，登绮席，泪珠滴，十二晚峰高历历。\par

\section{天仙子·踯躅花开红照水}
\textbf{皇甫松}\par\vspace{0.5em}
踯躅花开红照水，鹧鸪飞绕青山觜。\par
行人经岁始归来，千万里，错相倚，懊恼天仙应有以。\par

\section{浪淘沙·滩头细草接疏林}
\textbf{皇甫松}\par\vspace{0.5em}
滩头细草接疏林，浪恶罾船半欲沉。\par
宿鹭眠鸥飞旧浦，去年沙嘴是江心。\par

\section{浪淘沙·蛮歌豆蔻北人愁}
\textbf{皇甫松}\par\vspace{0.5em}
蛮歌豆蔻北人愁，蒲雨杉风野艇秋。\par
浪起鵁鶄眠不得，寒沙细细入江流。\par

\section{杨柳枝·春入行宫映翠微}
\textbf{皇甫松}\par\vspace{0.5em}
春入行宫映翠微，玄宗侍女舞烟丝。\par
如今柳向空城绿，玉笛何人更把吹。\par

\section{杨柳枝·烂漫春归水国时}
\textbf{皇甫松}\par\vspace{0.5em}
烂漫春归水国时，吴王宫殿柳丝垂。\par
黄莺长叫空闺畔，西子无因更得知。\par

\section{摘得新·酌一卮}
\textbf{皇甫松}\par\vspace{0.5em}
酌一卮，须教玉笛吹。\par
锦筵红蜡烛，莫来迟。\par
繁红一夜经风雨，是空枝。\par

\section{摘得新·摘得新}
\textbf{皇甫松}\par\vspace{0.5em}
摘得新，枝枝叶叶春。\par
管弦兼美酒，最关人。\par
平生都得几十度，展香茵。\par

\section{梦江南·兰烬落}
\textbf{皇甫松}\par\vspace{0.5em}
兰烬落，屏上暗红蕉。\par
闲梦江南梅熟日，夜船吹笛雨萧萧，人语驿边桥。\par

\section{梦江南·楼上寝}
\textbf{皇甫松}\par\vspace{0.5em}
楼上寝，残月下帘旌。\par
梦见秣陵惆怅事，桃花柳絮满江城，双髻坐吹笙。\par

\section{采莲子·菡萏香莲十顷陂}
\textbf{皇甫松}\par\vspace{0.5em}
菡萏香莲十顷陂（举棹），小姑贪戏采莲迟（年少）。\par
晚来弄水船头湿（举棹），更脱红裙裹鸭儿（年少）。\par

\section{采莲子·船动湖光滟滟秋}
\textbf{皇甫松}\par\vspace{0.5em}
船动湖光滟滟秋（举棹），贪看年少信船流（年少）。\par
无端隔水抛莲子（举棹），遥被人知半日羞（年少）。\par

\section{浣溪沙·清晓妆成寒食天}
\textbf{韦庄}\par\vspace{0.5em}
清晓妆成寒食天，柳球斜袅间花钿，卷帘直出画堂前。\par
指点牡丹初绽朵，日高犹自凭朱栏，含颦不语恨春残。\par

\section{浣溪沙·欲上秋千四体慵}
\textbf{韦庄}\par\vspace{0.5em}
欲上秋千四体慵，拟教人送又心忪。画堂帘幕月明风。\par
此夜有情谁不极，隔墙梨雪又玲珑，玉容憔悴惹微红。\par

\section{浣溪沙·惆怅梦余山月斜}
\textbf{韦庄}\par\vspace{0.5em}
惆怅梦余山月斜，孤灯照壁背窗纱，小楼高阁谢娘家。\par
暗想玉容何所似，一枝春雪冻梅花，满身香雾簇朝霞。\par

\section{浣溪沙·绿树藏莺莺玉啼}
\textbf{韦庄}\par\vspace{0.5em}
绿树藏莺莺玉啼，柳丝斜拂白铜堤，弄珠江上草萋萋。\par
日暮饮归何处客，绣鞍骢马一声嘶，满身兰麝醉如泥。\par

\section{浣溪沙·夜夜相思更漏残}
\textbf{韦庄}\par\vspace{0.5em}
夜夜相思更漏残，伤心明月凭栏干，想君思我锦衾寒。\par
咫尺画堂深似海，忆来唯把旧书看，几时携手入长安。\par

\section{菩萨蛮·红楼别夜堪惆怅}
\textbf{韦庄}\par\vspace{0.5em}
红楼别夜堪惆怅，香灯半卷流苏帐。\par
残月出门时，美人和泪辞。\par
琵琶金翠羽，弦上黄莺语。\par
劝我早归家，绿窗人似花。\par

\section{菩萨蛮·人人尽说江南好}
\textbf{韦庄}\par\vspace{0.5em}
人人尽说江南好，游人只合江南老。\par
春水碧于天，画船听雨眠。\par
垆边人似月，皓腕凝双雪。\par
未老莫还乡，还乡须断肠。\par

\section{菩萨蛮·如今却忆江南乐}
\textbf{韦庄}\par\vspace{0.5em}
如今却忆江南乐，当时年少春衫薄。\par
骑马倚斜桥，满楼红袖招。\par
翠屏金屈曲，醉入花丛宿。\par
此度见花枝，白头誓不归。\par

\section{菩萨蛮·劝君今夜须沉醉}
\textbf{韦庄}\par\vspace{0.5em}
劝君今夜须沉醉，樽前莫话明朝事。\par
珍重主人心，酒深情亦深。\par
须愁春漏短，莫诉金杯满。\par
遇酒且呵呵，人生能几何。\par

\section{菩萨蛮·洛陽城里春光好}
\textbf{韦庄}\par\vspace{0.5em}
洛陽城里春光好，洛陽才子他乡老。\par
柳暗魏王堤，此时心转迷。\par
桃花春水绿，水上鸳鸯浴。\par
凝恨对残晖，忆君君不知。\par

\section{归国遥·春欲暮}
\textbf{韦庄}\par\vspace{0.5em}
春欲暮，满地落花红带雨。\par
惆怅玉笼鹦鹉，单栖无伴侣。\par
南望去程何许，问花花不语。\par
早晚得同归去，恨无双翠羽。\par

\section{归国遥·金翡翠}
\textbf{韦庄}\par\vspace{0.5em}
金翡翠，为我南飞传我意。\par
罨画桥边春水，几年花下醉。\par
别后只知相愧，泪珠难远寄。\par
罗幕绣帷鸳被，旧欢如梦里。\par

\section{归国遥·春欲晚}
\textbf{韦庄}\par\vspace{0.5em}
春欲晚，戏蝶游蜂花烂漫。\par
日落谢家池馆，柳丝金缕断。\par
睡觉绿鬟风乱，画屏云雨散。\par
闲倚博山长叹，泪流沾皓腕。\par

\section{应天长·绿槐陰里黄莺语}
\textbf{韦庄}\par\vspace{0.5em}
绿槐陰里黄莺语，深院无人春昼午。\par
画帘垂，金凤舞，寂寞绣屏香一炷。\par
碧天云，无定处，空有梦魂来去。\par
夜夜绿窗风雨，断肠君信否。\par

\section{应天长·别来半岁音书绝}
\textbf{韦庄}\par\vspace{0.5em}
别来半岁音书绝，一寸离肠千万结。\par
难相见，易相别，又是玉楼花似雪。\par
暗相思，无处说，惆怅夜来烟月。\par
想得此时情切，泪沾红袖黦。\par

\section{荷叶杯·绝代佳人难得}
\textbf{韦庄}\par\vspace{0.5em}
绝代佳人难得，倾国，花下见无期。\par
一双愁黛远山眉，不忍更思惟。\par
闲掩翠屏金凤，残梦，罗幕画堂空。\par
碧天无路信难通，惆怅旧房栊。\par

\section{荷叶杯·记得那年花下}
\textbf{韦庄}\par\vspace{0.5em}
记得那年花下，深夜，初识谢娘时。\par
水堂西面画帘垂，携手暗相期。\par
惆怅晓莺残月，相别，从此隔音尘。\par
如今俱是异乡人，相见更无因。\par

\section{清平乐·春愁南陌}
\textbf{韦庄}\par\vspace{0.5em}
春愁南陌，故国音书隔。\par
细雨霏霏梨花白，燕拂画帘金额。\par
尽日相望王孙，尘满衣上泪痕。\par
谁向桥边吹笛，驻马西望销魂。\par

\section{清平乐·野花芳草}
\textbf{韦庄}\par\vspace{0.5em}
野花芳草，寂寞关山道。\par
柳吐金丝莺语早，惆怅香闺暗老。\par
罗带悔结同心，独凭朱栏思深。\par
梦觉半床斜月，小窗风触鸣琴。\par

\section{清平乐·何处游女}
\textbf{韦庄}\par\vspace{0.5em}
何处游女，蜀国多云雨。\par
云解有情花解语，窣地绣罗金缕。\par
妆成不整金钿，含羞待月秋千。\par
住在绿槐陰里，门临春水桥边。\par

\section{清平乐·莺啼残月}
\textbf{韦庄}\par\vspace{0.5em}
莺啼残月，绣阁香灯灭。\par
门外马嘶郎欲别，正是落花时节。\par
妆成不画蛾眉，含愁独倚金扉。\par
云路香尘莫扫，扫即郎去归迟。\par

\section{望远行·欲别无言倚画屏}
\textbf{韦庄}\par\vspace{0.5em}
欲别无言倚画屏，含恨暗伤情。\par
谢家庭树锦鸡鸣，残月落边城。\par
人欲别，马频嘶，绿槐千里长堤。\par
出门芳草路萋萋，云雨别来易东西。\par
不忍别君后，却入旧香闺。\par

\end{document}